\section{Empirical results}

\subsection{Data-dependent transition in training-set reproduction}

% AUTHOR NOTE: Present UNet and DiT novelty curves with uncertainty and exact
% update budgets. Separate equal-budget comparisons from longer-training tests.
% State where behavior is reproducible and where it remains anomalous.
\draftplaceholder

% AUTHOR NOTE: Insert a figure showing the nearest-training transition and
% representative generated, nearest-training, and absolute-difference panels.

\subsection{Novelty and statistical validity do not coincide}

% AUTHOR NOTE: Compare novelty with one-point and power-spectrum errors across
% training-set size. Quantify any region in which novelty passes while physical
% statistics do not, including seed or sample uncertainty before naming it a
% systematic regime.
\draftplaceholder

% AUTHOR NOTE: Insert the full data-size sweep of one-point statistics and
% scale-resolved power-spectrum ratios or errors.

\subsection{Conditional parameter recovery}

% AUTHOR NOTE: Compare requested and recovered cosmological parameters for
% selected memorizing and generalizing conditional models. Include held-out-real
% probe validation and posterior coverage before interpreting confidence.
\draftplaceholder

\subsection{Architecture and optimization-time dependence}

% AUTHOR NOTE: Analyze checkpoint trajectories, loss curves, sampler controls,
% and architecture comparisons. State whether additional training changes
% novelty, physical statistics, both, or neither. Do not infer convergence from
% denoising loss alone.
\draftplaceholder
